% 08-evaluacion-resultados.tex — Fuente: docs/informe-final/08-evaluacion-resultados.md

\chapter{Evaluación empírica y resultados}
\label{cap:evaluacion-resultados}

Cada bloque se reporta con: (i) pregunta de investigación asociada, (ii)
métrica, (iii) datos crudos y su ubicación, (iv) estadística descriptiva,
y (v) interpretación textual.

\section{Rendimiento (k6) --- RQ1}

\textbf{Métrica:} percentiles de latencia (\texttt{http\_req\_duration}) y
tasa de error HTTP $\geq$500 sobre
\texttt{GET /api/v1/catalogo?page=0\&size=20}, 50~VUs, 30~s, 3 corridas
por escenario. \textbf{Datos crudos:}
\texttt{docs/mediciones/perf/k6-run\{1,2,3\}.json} (caliente),
\texttt{k6-cold-run\{1,2,3\}.json} (frío); consolas en
\texttt{k6-console-run*.txt}.

\begin{table}[H]
\centering
\caption{Resultados de la prueba de carga k6 por escenario}
\label{tab:k6}
\resizebox{\textwidth}{!}{%
\begin{tabular}{@{}lrrrrrlr@{}}
\toprule
\textbf{Escenario} & \textbf{n} & \textbf{Media (ms)} & \textbf{Mediana (ms)} & \textbf{DT (ms)} & \textbf{IC 95\,\% (ms)} & \textbf{p50/p90/p95/p99 (ms)} & \textbf{Error $\geq$500} \\
\midrule
Caliente & 4500 & 21.38 & 13.01 & 32.71 & [20.42, 22.33] & 13.01 / 31.75 / \textbf{50.17} / 205.60 & 0.00\,\% \\
Frío     & 4500 & 13.62 &  9.04 & 14.52 & [13.20, 14.05] &  9.04 / 22.39 / \textbf{39.14} /  89.98 & 0.00\,\% \\
\bottomrule
\end{tabular}}
\end{table}

\textbf{Umbral:} p95 $\leq$200~ms caliente (\textbf{cumple}, 50.17~ms),
p95 $\leq$500~ms frío (\textbf{cumple}, 39.14~ms). Throughput
$\approx$48.5--49.0~req/s por corrida.

\textbf{Interpretación:} ambos escenarios cumplen holgadamente sus
umbrales respectivos. El hallazgo más relevante no es el cumplimiento en
sí, sino que el escenario ``frío'' resultó \textbf{más rápido} que el
``caliente'' --- contraintuitivo, y documentado con honestidad
metodológica en \texttt{docs/mediciones/perf/REPORTE-PERF.md}: el script
de carga golpea siempre la misma combinación de parámetros de consulta,
por lo que la clave de caché de Spring (\texttt{@Cacheable},
Capítulo~\ref{cap:implementacion}~\S7.2) se puebla con la primera
petición en ambos escenarios --- de las 1500 iteraciones por corrida, como
mucho una es un \emph{miss} real contra PostgreSQL. La diferencia de
medias observada (21.38~ms vs. 13.62~ms) es ruido de ejecución (JIT/GC de
la JVM, carga del contenedor), no un efecto real de caché frío vs.
caliente. \textbf{No se ejecutó} el test inferencial (Wilcoxon) ni el
tamaño de efecto que exige el Bloque~C de la guía para esta comparación
--- protocolo recomendado por Arcuri y Briand~\cite{ARCURI2011},
declarado como brecha en el Capítulo~\ref{cap:materiales-metodos}~\S5.6 y
retomado como amenaza a la validez en el Capítulo~\ref{cap:validez}.

\section{Seguridad (OWASP + ZAP + análisis estático) --- RQ1, RQ2}

\textbf{Métrica:} estado (aprobado/hallazgo) de 6 controles
OWASP~\cite{OWASP2021} evidenciados manualmente, más el resultado
agregado del escaneo automatizado ZAP baseline y del análisis estático
SpotBugs + find-sec-bugs~\cite{HALFOND2005,GOULD2004}. \textbf{Datos
crudos:} \texttt{docs/mediciones/sec/owasp/*.txt}, \texttt{sec/zap/},
\texttt{sec/static-analysis/}.

\begin{table}[H]
\centering
\caption{Controles OWASP evidenciados manualmente}
\label{tab:owasp}
\begin{tabular}{@{}p{5.4cm}p{2.0cm}p{5.6cm}@{}}
\toprule
\textbf{Control OWASP} & \textbf{Estado} & \textbf{Evidencia} \\
\midrule
A01 --- Control de acceso roto & Aprobado & \texttt{sec/owasp/a01-control-acceso.txt} \\
A02 --- Fallas criptográficas (TLS) & Aprobado & \texttt{sec/owasp/a02-tls.txt} \\
A03 --- Inyección & Aprobado & \texttt{sec/owasp/a03-inyeccion.txt} \\
A05 --- Configuración de seguridad incorrecta & Aprobado & \texttt{sec/owasp/a05-cabeceras.txt} \\
A07 --- Fallas de identificación y autenticación & Aprobado & \texttt{sec/owasp/a07-rate-limit.txt} \\
A09 --- Fallas de registro y monitoreo & Aprobado & \texttt{sec/owasp/a09-logging.txt} \\
\bottomrule
\end{tabular}
\end{table}

A04, A06, A08, A10 quedan fuera de alcance declarado (no evaluados) --- ver
\texttt{docs/mediciones/sec/REPORTE-SEC.md}.

\textbf{ZAP baseline:} \texttt{FAIL-NEW: 0} $\cdot$ \texttt{WARN-NEW: 8}
(severidad media/baja, ninguno bloqueante) $\cdot$ \texttt{PASS: 59}.
Cumple el umbral de ``sin hallazgos altos''.

\textbf{Análisis estático (SpotBugs + find-sec-bugs):} \textbf{0
hallazgos} de inyección SQL
(\texttt{SQL\_NONCONSTANT\_STRING\_PASSED\_TO\_EXECUTE},
\texttt{SQL\_INJECTION\_JDBC} y variantes) sobre el bytecode del backend
completo --- consistente con el modelo de detección estático propuesto
originalmente por Gould, Su y Devanbu~\cite{GOULD2004}.

\textbf{Interpretación:} la combinación de evidencia manual reproducible
(6 controles) y escaneo automatizado (ZAP + análisis estático) da una
cobertura de seguridad consistente en ambas direcciones (top-down por
control OWASP, bottom-up por herramienta automatizada), sin
contradicciones entre ambas fuentes.

\section{Usabilidad (SUS) --- RQ1}

\textbf{Métrica:} puntaje System Usability Scale
(0--100)~\cite{BROOKE1996}, N=16 (umbral mínimo de la guía: N$\geq$15).
\textbf{Datos crudos:} \texttt{docs/mediciones/sus/sus-raw.csv}, perfil de
participantes en \texttt{perfil-participantes.csv}.

\begin{table}[H]
\centering
\caption{Resultados del cuestionario SUS}
\label{tab:sus}
\begin{tabular}{@{}lr@{}}
\toprule
\textbf{Métrica} & \textbf{Valor} \\
\midrule
n & 16 \\
Media & 76.88 \\
Mediana & 77.50 \\
DT & 14.48 \\
IC 95\,\% & [69.16, 84.59] \\
Interpretación (escala de Bangor~\cite{BANGOR2009}) & \textbf{B} \\
Umbral del proyecto ($>$68) & Supera \\
\bottomrule
\end{tabular}
\end{table}

\begin{figure}[H]
\centering
\includegraphics[width=0.6\textwidth]{boxplot-sus.png}
\caption{Diagrama de caja de los puntajes SUS ($N=16$). Fuente:
\texttt{docs/mediciones/sus/boxplot-sus.png}.}
\label{fig:boxplot-sus}
\end{figure}

\textbf{Nota de conformidad:} el Bloque~C de la guía exige gráficos
vectoriales (SVG/PDF); la Figura~\ref{fig:boxplot-sus} está en formato
raster (PNG). Regenerarlo en SVG desde el mismo script
(\texttt{docs/mediciones/sus/graficar-sus.py}) es una corrección de bajo
esfuerzo pendiente antes del cierre.

\textbf{Interpretación:} el puntaje medio (76.88) está por encima del
percentil~50 de la escala SUS normalizada y en la categoría~B de Bangor
(rango ``excelente'' bajo), superando el umbral mínimo del proyecto
($>$68) con margen. La limitación de muestreo por conveniencia
(Capítulo~\ref{cap:materiales-metodos}~\S5.4) restringe la generalización
de este resultado a la población real de usuarios --- retomado en el
Capítulo~\ref{cap:validez}.

\section{Cobertura de código (JaCoCo) --- RQ1}

\textbf{Métrica:} cobertura de líneas, ramas y complejidad ciclomática
sobre los módulos de dominio, servicios y controladores --- las mismas
tres dimensiones que Zhu, Hall y May~\cite{ZHU1997} sistematizan como
criterios de adecuación de pruebas. \textbf{Datos crudos:}
\texttt{docs/mediciones/jacoco/report.xml}, \texttt{html/jacoco.csv}.

\begin{table}[H]
\centering
\caption{Cobertura de código JaCoCo}
\label{tab:jacoco}
\begin{tabular}{@{}lllc@{}}
\toprule
\textbf{Métrica} & \textbf{Cobertura} & \textbf{Umbral} & \textbf{Cumple} \\
\midrule
Lines         & 72.0\,\% (2867/3981)   & $\geq$70\,\% & \checkmark \\
Branches      & 62.5\,\% (703/1125)    & $\geq$70\,\% & $\times$ \\
Instructions  & 69.7\,\% (12825/18404) & ---          & --- \\
Complexity    & 56.5\,\% (775/1371)    & ---          & --- \\
\bottomrule
\end{tabular}
\end{table}

\textbf{Interpretación:} la cobertura de líneas cumple el umbral de
70\,\%; \textbf{la cobertura de ramas no lo alcanza} (62.5\,\% vs.
70\,\%) --- declarado sin maquillaje, es una brecha real, no un redondeo
favorable. Cabe además la advertencia empírica de Inozemtseva y
Holmes~\cite{INOZEMTSEVA2014}: un porcentaje de cobertura alto no implica
por sí solo una suite de pruebas efectiva para detectar defectos, por lo
que estas cifras deben leerse como una medida de exhaustividad estructural,
no como un sustituto de la efectividad real de las 522 pruebas. La
medición es un agregado global del backend; la guía exige desglose por
módulo (dominio, servicios, controladores) que
\texttt{docs/mediciones/jacoco/REPORTE-JACOCO.md} no provee todavía de
forma separada --- pendiente para el cierre. Además, \texttt{pom.xml}
configura JaCoCo solo para reportar, sin \texttt{jacoco:check} que falle
el build bajo el umbral (Capítulo~\ref{cap:diseno-arquitectura}, brecha
declarada también en \texttt{docs/observaciones/INFORME-BRECHAS-ENTREGA-FINAL.md}).

\section{Calidad web (Lighthouse) --- RQ1}

\textbf{Métrica:} puntajes Performance, Accessibility, Best Practices, SEO
(0--100), 3 corridas por perfil (mobile, desktop). \textbf{Datos crudos:}
\texttt{docs/mediciones/lighthouse/lhci-20260817-*-\{mobile,desktop\}-run\{1,2,3\}.json}.

\begin{table}[H]
\centering
\caption{Resultados Lighthouse por perfil y categoría}
\label{tab:lighthouse}
\begin{tabular}{@{}lllcc@{}}
\toprule
\textbf{Categoría} & \textbf{Mobile (3 corridas)} & \textbf{Desktop (3 corridas)} & \textbf{Umbral} & \textbf{Cumple} \\
\midrule
Performance     & 81 / 81 / 80   & 100 / 100 / 100 & $\geq$80 & \checkmark ambos \\
Accessibility   & 93 / 93 / 93   & 93 / 93 / 93    & $\geq$90 & \checkmark ambos \\
Best Practices  & 96 / 96 / 96   & 96 / 96 / 96    & $\geq$90 & \checkmark ambos \\
SEO             & 100 / 100 / 100 & 100 / 100 / 100 & $\geq$90 & \checkmark ambos \\
\bottomrule
\end{tabular}
\end{table}

\textbf{Interpretación:} las cuatro categorías cumplen umbral en ambos
perfiles, con consistencia prácticamente perfecta entre las 3 corridas de
cada perfil (varianza $\leq$1 punto en Performance mobile, 0 en el
resto). Los hallazgos reales que impiden el 100/100 están documentados
sin maquillaje en \texttt{docs/mediciones/lighthouse/REPORTE-LIGHTHOUSE.md}:
\texttt{target-size} (elementos táctiles bajo el tamaño mínimo
recomendado, Accessibility 93) y un \texttt{403} de red esperado en
\texttt{/api/auth/refresh} al cargar sin sesión (Best Practices 96) ---
ninguno de los dos es un defecto real de la aplicación, sino
comportamiento esperado bajo las condiciones de la auditoría.

\section{Resumen de reproducibilidad de las figuras y tablas de este capítulo}

Todo dato reportado en este capítulo proviene de un archivo crudo
versionado bajo \texttt{docs/mediciones/} y de un script o comando
documentado (ver Capítulo~\ref{cap:materiales-metodos}~\S5.5 y
\texttt{docs/mediciones/DATA-PROVENANCE.md} para el commit/script exacto
de cada medición) --- ningún número de esta sección se calculó fuera del
repositorio.
